\subsection{Overview}

The goal of our experiments is to exemplify the utility of the Before-ARC generator to study compositionality and generalization. Furthermore, we aim to demonstrated how it enables targeted evaluation of these two abilities, in ways that other similar data environments do not. 

When possible, we attempted to align the experiment settings accross different data environments. We however argue that there are several experiment settings of strong scientific relevance to the study of compositionality and generalization are made possible in Before-ARC, but not in the other data environments.

\subsubsection{Before-ARC configuration}

The Before-ARC data environment follows a \textit{generation} output mode - given an input grid, models must generate the corresponding output grid. In the \textbf{Generalization} study, the task (i.e. the sequence of transformations) is kept unchanged for the given experiment setting, while other parameters such as the grid size and number of shapes per grid varies. The model must then learn to correctly predict the output grids given grid or shapes configurations that differ between train and test time. In the \textbf{Compositionality} study, the model must learn to apply a sequence of transformations to a grid. In some experiment settings, there are multiple transformation sequence that the model must learn to execute at train time. The transformation sequences are also always different between training and testing. As such, we must provide some context to inform the model on which task it must perform. We do so by adding a task embedding, similarly to CVR \cite{cvr2022}. The task embedding is a sequence of tokens that indicates to the model the transformations it must apply (in the correct order). The task embedding is simply appended to the end of the input sequence. 
 
\subsubsection{CVR Configuration}

For the CVR \cite{cvr2022} data environment, consider a \textit{classification} output mode, which the authors call "odd-one-out". In this framework, a sample contrains $k=4$ RGB input images following an AVR rule, and the output is a categorical value indicating which of the $k$ images is the odd one. In the \textbf{Generalization} study, we use a function built in to CVR's code architecture which simultaneously varies different parameters such as the number, size and complexity of the shape. We thus train our models on a given sample "difficulty" and test on a bigger difficulty. In the \textbf{Compositionality} study, we consider the set of 103 tasks that they offer. There are 9 elementary rules and 94 composite ones. The composite rules are either roughly the result of a combination (e.g., \texttt{task\_pos\_contact or task\_size\_count\_2}), or they are a "standalone" more complex rule (e.g., \texttt{task\_sym\_mir}). In our compositionality study, due to resource constraints, we focus on a subset of 4 elementary and composite rules built out these 4 elementary ones. Similarly to the Before-ARC setting, we build experiment settings in order to train on a given set of rules, and test on a different one - similarly to the Before-ARC setting.

\subsubsection{RE-ARC Configuration}\label{rearc}

The RE-ARC \cite{reverse_engineer_hodel} data environment follows a \textit{generation} output mode, similarly to the Before-ARC data. In this data environment, we only consider the generalization study, as there is no systematic way to perform compositionality study. Specifically, for the \textbf{Generalization} study, we randomly sample 10 tasks out of the 400 possible tasks, and conduct our various experiments on those. RE-ARC allows for a difficulty parameter (with lower and upper bound) to be set. As such, we train on a specific difficulty interval ([0-0.7]) and test on a higher difficulty interval ([0.7-1]). The difficulty interval will be characterized differently on each tasks: each task allows for different variations that can modify its complexity. As such, the exact difficulty range of the task is not the same across tasks, which makes the study of generalization less precise.
